\documentclass[10pt,a4paper]{article}
\usepackage[utf8]{inputenc}
\usepackage[T1]{fontenc}
\usepackage[slovene]{babel}
\usepackage{lmodern}
\usepackage[a4paper,margin=1.9cm,top=1.6cm,bottom=1.8cm]{geometry}
\usepackage{graphicx}
\usepackage{booktabs}
\usepackage{float}
\usepackage{xcolor}
\usepackage[hidelinks]{hyperref}
\usepackage{xurl}
\usepackage[font=small,labelfont=bf,skip=3pt]{caption}
\usepackage{parskip}
\setlength{\parskip}{4pt}
\usepackage{titlesec}
\titlespacing*{\section}{0pt}{8pt}{3pt}
\titlespacing*{\subsection}{0pt}{6pt}{2pt}
\graphicspath{{../figures/a/}{../figures/b/}{../figures/c/}}
\newcommand{\m}{$-$}

\begin{document}

\begin{center}
{\Large\bfseries Neto pozicija madžarskega elektroenergetskega sistema}\\[3pt]
Avtor: Veronika Tomažič
\end{center}

\begin{itemize}\setlength\itemsep{2pt}
\item \textbf{a)} Madžarska elektriko stalno uvaža: od leta 2015 v povprečju 1000–1400\,MW, kar je približno
četrtina porabe. Skoraj polovico domače elektrike proizvede jedrska elektrarna Paks, ki dela ves čas enako.
Plinske elektrarne delajo, ko je elektrika draga, sončnih elektrarn pa je vsako leto več in opoldne uvoz
skoraj izničijo. Elektriko kupuje na severu (Slovaška, Avstrija), del jo proda na jug.
\item \textbf{b)} Za teden 47 napovedujemo povprečni neto položaj \textbf{\m1873\,MW} (uvoz, $\approx$\,\m315\,GWh),
z razponom \textbf{\m2146 do \m1483\,MW} (80\,\%). Napoved je izbrana med 18 modeli, preverjenimi na 202 tednih.
\item \textbf{c)} Povprečna dan-vnaprej cena v tednu 47 bo po naši oceni \textbf{nižja} od tedna 45
(290,98\,€/MWh), približno \textbf{260–275\,€/MWh}. Glavni razlog je precej močnejši veter v Nemčiji, ki
poceni elektriko v celi regiji, medtem ko je cena plina enaka.
\end{itemize}

Analiza temelji na javnih podatkih od leta 2015 do 15.\,11.\,2026 v 15-minutni ločljivosti: proizvodnja po
vrstah elektrarn, poraba, čezmejno trgovanje in dan-vnaprej cene (Energy-Charts, ki povzema uradne podatke
ENTSO-E), vreme za pet večjih madžarskih mest in veter v severni Nemčiji (Open-Meteo) ter cena plina TTF
(investing.com). Vsi časi so v srednjeevropskem času (CET).

% =====================================================================
\subsection*{a) Neto pozicija Madžarske}

Neto pozicija je razlika med proizvodnjo in porabo in je enaka razliki med izvozom in uvozom.
Predznak je kot v nalogi: \textbf{+ pomeni izvoz, \m\ pomeni uvoz}.

\subsection*{Kaj proizvaja elektriko}

Madžarska ima nekaj glavnih vrst elektrarn, ki delujejo zelo različno (tabela~\ref{tab:teh}). Jedrska
elektrarna Paks proizvede skoraj polovico domače elektrike in dela ves čas z enako močjo. Plinske elektrarne
delajo takrat, ko je elektrika draga, na primer zvečer in pozimi. Ko je elektrika poceni, jih ugasnejo in
elektriko raje uvozijo, zato se njihova proizvodnja najbolj spreminja. Sončne elektrarne proizvajajo samo
podnevi in ob lepem vremenu, njihova elektrika pa je skoraj zastonj. Lignitna elektrarna Mátra se postopno
zapira, vetrnih in hidroelektrarn pa je na Madžarskem malo.

\begin{table}[H]
\centering\small
\footnotesize\setlength{\tabcolsep}{4pt}
\begin{tabular}{lrrll}
\toprule
vrsta & TWh & delež & način delovanja & strošek / vloga pri ceni \\
\midrule
jedrska (Paks) & 15,2 & 45\,\% & enakomerno; izpadi ob revizijah & nizek \\
plin & 7,4 & 22\,\% & prilagodljivo, ob visoki ceni & visok; pogosto določa ceno \\
sonce & 6,8 & 20\,\% & samo podnevi, odvisno od vremena & skoraj nič; opoldne zniža ceno \\
lignit & 1,5 & 5\,\% & stabilno, upada & srednji \\
biomasa, veter, voda, drugo & 2,9 & 8\,\% & majhni in stabilni & nizek \\
\midrule
\textbf{proizvodnja / poraba / neto} & \multicolumn{4}{l}{\textbf{34,0 / 43,2 / \m9,5 TWh} (uvoz $\approx$ 22\,\% porabe)} \\
\bottomrule
\end{tabular}
\caption{Vrste elektrarn (vse elektrarne iste vrste skupaj), leto 2025.}
\label{tab:teh}
\end{table}

\begin{figure}[H]
\centering
\includegraphics[width=0.6\linewidth]{a1b_letni_trend_s_soncem.pdf}
\caption{Proizvodnja po vrstah elektrarn in poraba, letna povprečja 2020–2026.}
\label{fig:trend}
\end{figure}

Slika~\ref{fig:trend} kaže, kako se je to spreminjalo od leta 2020; prej v podatkih ni sončnih elektrarn,
zato starejših let na sliki ni. Razlika med stolpci in črto porabe je uvoz. Lignita je vsako leto manj (s 681\,MW
leta 2015 na 160\,MW leta 2026), sončnih elektrarn pa vsako leto več (s 181\,MW leta 2020 na 946\,MW leta 2026).
Letos proizvedejo skoraj toliko kot plin. Paks je bil 10 let skoraj enak ($\sim$1720\,MW), letos pa je nižji
(1558\,MW), ker je poleti zaradi nizke Donave delal manj. Uvoz je ves čas ostal podoben, med 1000 in 1400\,MW
(tudi v letih 2015–2019), ker je izpad lignita in Paksa nadomestilo sonce.


\subsection*{Kaj poganja porabo}

Poraba je skozi leta skoraj enaka, v povprečju $\sim$4900\,MW, iz dneva v dan pa se precej spreminja. Odvisna je
predvsem od temperature in od tega, ali je delovni dan ali vikend.

Ko je toplo (nad $\sim$15\,°C), temperatura na porabo skoraj ne vpliva. Ko je hladneje, ljudje ogrevajo, zato
vsaka nižja stopinja poveča porabo za $\sim$95\,MW (slika~\ref{fig:temp}, desni graf: 2022–2026). Poleti, ko je
zelo vroče (nad $\sim$22\,°C), porabo spet poveča hlajenje. Če bi vpliv temperature ocenili z eno ravno črto čez
vse leto, bi dobili le $\sim$47\,MW na stopinjo in bi vpliv mraza podcenili.

Slika~\ref{fig:temp} kaže tri obdobja: vsa leta (2015–2026), leta s podatki o sončnih elektrarnah (2020–2026) in
leta z najboljšimi podatki (2022–2026). Vzorec je v vseh enak, vpliv mraza pa se povečuje (s $\sim$66 na
$\sim$94\,MW na stopinjo), verjetno zaradi vse več toplotnih črpalk.

Ob vikendih je poraba nižja, ker tovarne in podjetja ne delajo: v soboto za $\sim$9\,\%, v nedeljo za $\sim$13\,\%.

\begin{figure}[H]
\centering
\includegraphics[width=0.8\linewidth]{a5a_poraba_temperatura.pdf}
\caption{Dnevna poraba proti dnevni temperaturi (delovni dnevi brez praznikov) za tri obdobja. Ravna črta
(črtkano) proti zalomljeni črti, ki loči ogrevanje, nevtralno območje in hlajenje.}
\label{fig:temp}
\end{figure}

\subsection*{Prek katerih meja trguje}

Madžarska je povezana z vsemi sedmimi sosednjimi državami. V tabeli~\ref{tab:meje} je za vsako mejo povprečje
trgovanja v letu (+ izvoz, \m\ uvoz). Slovaška in Avstrija imata vsa leta minus: od tam Madžarska elektriko kupuje,
iz Slovaške največ ($\sim$1000–1250\,MW). Hrvaška ima vsa leta plus: tja Madžarska elektriko prodaja. Pri Romuniji,
Srbiji in Sloveniji se smer menja iz leta v leto. Madžarska torej elektriko kupuje na severu, del je porabi doma,
del pa proda naprej na jug.

Pri pregledu podatkov sem opazila, da številke po posameznih mejah niso posli med dvema državama, ampak jih
izračuna algoritem skupnega evropskega trga. Najbolj zanesljiva je zato njihova vsota, to je neto pozicija.

\begin{table}[H]
\centering\small


\begin{tabular}{lrrrrrrrr}
\toprule
 & Slovaška & Avstrija & Romunija & Hrvaška & Srbija & Slovenija & Ukrajina & \textbf{skupaj} \\
\midrule
2024 & \m1029 & \m893 & +378 & +281 & \m38 & \m153 & +185 & \textbf{\m1268} \\
2025 & \m1246 & \m746 & +612 & +350 & \m82 & \m34 & +65 & \textbf{\m1082} \\
2026 & \m1099 & \m495 & \m65 & +335 & \m195 & +102 & +186 & \textbf{\m1233} \\
\bottomrule
\end{tabular}
\caption{Neto trgovanje po mejah, MW (letna povprečja; + izvoz iz HU, \m\ uvoz v HU). 2026: do 15.\,11.}
\label{tab:meje}
\end{table}

\subsection*{Kako se spreminja čez dan in čez leto}

\textbf{Čez leto:} pozimi Madžarska uvozi približno dvakrat več elektrike kot spomladi in poleti. Poleti uvoz
zmanjšajo sončne elektrarne, pozimi pa sonca skoraj ni in je poraba zaradi ogrevanja večja, zato več delajo
plinske elektrarne, ostalo pa se uvozi.

\textbf{Čez dan:} ker napovedujemo teden v novembru, smo dnevni potek pogledali za oktober in november v
preteklih letih (slika~\ref{fig:profil}). Opoldne, ko proizvaja sonce, je uvoz danes skoraj nič, leta 2020 pa je
bil še $\sim$1300\,MW. Zvečer, ko sonca ni in je poraba največja, Madžarska uvozi največ, $\sim$2300–2500\,MW.
Ta vzorec se ponavlja vsak dan in ga vidimo tudi v napovedi za teden 47.

\begin{figure}[H]
\centering
\includegraphics[width=0.6\linewidth]{a4b_profil_okt_nov_po_letih.pdf}
\caption{Neto pozicija čez dan v oktobru in novembru, po letih (15-min vrednosti, lokalni čas).
Opoldne sonce uvoz skoraj izniči, zvečer je uvoz največji.}
\label{fig:profil}
\end{figure}

% =====================================================================
\subsection*{b) Napoved neto položaja za teden 47}

\subsection*{Kaj bo vplivalo na neto položaj}

Napoved temelji na formuli iz naloge: \emph{neto = proizvodnja \m\ poraba} (model po formuli, na
sliki~\ref{fig:modeli} označen z 10b). Zato sem najprej pogledala, kateri deli formule se iz tedna v teden
najbolj spreminjajo, saj za napoved ni pomembno, kdo proizvede največ, ampak kdo se najbolj spreminja. Najbolj se
spreminjajo poraba, plinske elektrarne, sonce in Paks, ostale elektrarne (veter, lignit, biomasa, voda) pa skoraj
nič. Za teden 47 to pomeni:
\begin{itemize}\setlength\itemsep{1pt}
\item \textbf{temperatura → poraba:} v tednu 47 je napovedana povprečna temperatura 7,2\,°C (teden 46: 6,5\,°C);
začetek tedna je toplejši (do 10,5\,°C v sredo), konec tedna hladnejši (3,5–4,7\,°C);
\item \textbf{sonce:} napovedano je več sonca kot v oblačnem tednu 46 (sevanje $\sim$1,3\,kWh/m$^2$ na dan proti
$\sim$0,9), a manj kot v tednu 45;
\item \textbf{Paks:} en reaktor je do konca decembra na rednem vzdrževanju (uradno obvestilo na ENTSO-E), zato bo
Paks proizvajal enako kot v tednu 46, $\sim$1460\,MW;
\item \textbf{plin:} najbolj negotov, ker plinske elektrarne delajo samo, ko je elektrika draga. Ko je elektrika v
regiji poceni, ugasnejo, Madžarska pa več elektrike uvozi iz sosednjih držav.
\end{itemize}

\subsection*{Pristop: 18 modelov, preverjenih na 202 tednih}

Ker je odločitev za en sam pristop lahko naključna, smo primerjali 18 modelov, od najpreprostejših do
nevronskih mrež (slika~\ref{fig:modeli}). Vsak je moral za vsak teden od januarja 2023 do 15.\,11.\,2026
(202 tedna) napovedati naslednji teden \emph{samo s podatki do nedelje pred njim}, njegovo napako pa smo
primerjali z naivno napovedjo »enako kot prejšnji teden« (ki ima vrednost 1). Skupine modelov:
\begin{itemize}\setlength\itemsep{1pt}
\item \textbf{preprosti}: kot prejšnji teden, kot lani, povprečje zadnjih 2 ali 4 tednov;
\item \textbf{en dejavnik}: zadnja dva tedna + popravek samo za temperaturo, samo za sonce, samo za Paks ali samo za plin;
\item \textbf{vse skupaj}: linearna regresija z vsemi dejavniki, model »po formuli« (vsak vir posebej, nato
proizvodnja \m\ poraba), podobni dnevi iz preteklosti, nevronski mreži (dnevna in za ves 15-min profil);
\item \textbf{kombinacije}: povprečje več modelov.
\end{itemize}

\begin{figure}[H]
\centering
\includegraphics[width=0.8\linewidth]{b3_primerjava_modelov.pdf}
\caption{Napaka tedenske napovedi glede na naivni model (manj je bolje), test 202 tednov 2023–2026.}
\label{fig:modeli}
\end{figure}

Ugotovitve testa: najboljša sta linearna regresija in kombinacija regresije, modela po formuli in podobnih
dni – oba za \textbf{34\,\% boljša od naivne napovedi}; model po formuli iz naloge je takoj za njima (31\,\%).
Izmed posameznih dejavnikov najbolj pomaga sonce, plin sam pa nič, ker je preveč nepredvidljiv.
Nevronski mreži nista boljši od regresije – za zapletene modele imamo premalo podatkov ($\sim$1300 dni), zato se
lahko naučijo tudi šuma. Kot končni model smo izbrali \textbf{kombinacijo}, ker je bila v jesenskih tednih
(oktober–december) najboljša. Pravilo izbire je bilo določeno vnaprej, pred napovedjo tedna 47.
Opomba: v testu je uporabljeno izmerjeno vreme, zato je test nekoliko preveč optimističen.

\subsection*{Napoved}

Napoved je narejena z izbranim modelom samo s podatki do nedelje 15.\,11.\,2026, 23:45. Tabela~\ref{tab:napoved}
kaže napoved po dnevih. Stolpci pomenijo:
\begin{itemize}\setlength\itemsep{1pt}
\item \textbf{povprečna moč (MW)}: koliko elektrike Madžarska v povprečju uvaža v vsakem trenutku dneva
(povprečje 96 četrtur dneva; \m\ pomeni uvoz);
\item \textbf{energija (GWh)}: koliko elektrike skupaj uvozi v tem dnevu (povprečna moč × 24\,h);
\item \textbf{največ / najmanj uvoza}: najmočnejša in najšibkejša četrtura dneva.
\end{itemize}
\textbf{Teden} pri moči je povprečje sedmih dni, pri energiji pa njihova vsota; povezuje ju število ur v tednu: $314{,}7\,\mathrm{GWh} / 168\,\mathrm{h} = 1873\,\mathrm{MW}$.

\begin{table}[H]
\centering\small


\begin{tabular}{lrrll}
\toprule
dan & povprečna moč (MW) & energija (GWh) & največ uvoza & najmanj uvoza \\
\midrule
pon 16.\,11. & \m2201 & \m52,8 & \m3121 ob 19:15 & \m925 ob 11:15 \\
tor 17.\,11. & \m1942 & \m46,6 & \m2862 ob 19:15 & \m666 ob 11:15 \\
sre 18.\,11. & \m1732 & \m41,6 & \m2652 ob 19:15 & \m457 ob 11:15 \\
čet 19.\,11. & \m1811 & \m43,5 & \m2731 ob 19:15 & \m536 ob 11:15 \\
pet 20.\,11. & \m1890 & \m45,4 & \m2810 ob 19:15 & \m614 ob 11:15 \\
sob 21.\,11. & \m1790 & \m43,0 & \m2710 ob 19:15 & \m515 ob 11:15 \\
ned 22.\,11. & \m1748 & \m41,9 & \m2668 ob 19:15 & \m472 ob 11:15 \\
\midrule
\textbf{teden 47} & \textbf{\m1873} (povprečje) & \textbf{\m314,7} (vsota) & & \\
razpon (80\,\%) & \m2146 do \m1483 & \m360,5 do \m249,1 & & \\
\bottomrule
\end{tabular}
\caption{Napoved neto položaja za teden 47 po dnevih (\m\ = uvoz; lokalni čas CET).}
\label{tab:napoved}
\end{table}

\begin{figure}[H]
\centering
\includegraphics[width=0.85\linewidth]{b1_napoved_t47_izbrani.pdf}
\caption{Napoved neto pozicije za teden 47 po 15 minutah (lokalni čas) z razponom 10–90\,\%.
Črtkana vodoravna črta je povprečje tedna (\m1873\,MW).}
\label{fig:napoved}
\end{figure}

Razpon je izračunan iz napak modela v jesenskih tednih testa: v 8 od 10 podobnih tednov je bil dejanski
neto položaj v tem pasu. Vsak dan je uvoz najmanjši okoli 11:00, ko proizvaja sonce, in največji okoli 19:00,
ko sonca ni in je poraba visoka. Največ uvoza pričakujemo v ponedeljek (hladneje in oblačno), najmanj v
sredo (najtopleje, več sonca). Ob vikendu je poraba nižja, a je hladneje, zato se uvoz ne zmanjša veliko.

% =====================================================================
\subsection*{c) Ali bo cena v tednu 47 višja ali nižja od 290,98\,€/MWh?}

\subsection*{Logika: merit order}

Elektrarne so razvrščene od najcenejše do najdražje (sonce, veter, voda → jedrska → lignit → plin), ceno pa
določa najdražja, ki je še potrebna. Za Madžarsko so poleg domačih elektrarn del te »vrste« tudi
sosednji trgi: ko je elektrika pri sosedih poceni, Madžarska uvaža namesto da bi vklopila plin. Zato je cena
približno \textbf{cena plina × faktor}: faktor je večji, ko je doma veliko »preostale porabe« (poraba \m\ sonce
\m\ veter \m\ Paks – to, kar morata pokriti plin in uvoz), in manjši, ko v regiji piha
(slika~\ref{fig:merit}). Razmerje med ceno elektrike in ceno plina TTF je bilo v letih 2021–2026 večinoma
med 2 in 3.

\begin{figure}[H]
\centering
\includegraphics[width=0.6\linewidth]{c1_merit_order.pdf}
\caption{Razmerje med ceno elektrike in ceno plina proti preostali porabi na Madžarskem (dnevne vrednosti,
zadnjih 12 mesecev). Več preostale porabe pomeni dražje elektrarne in višjo ceno.}
\label{fig:merit}
\end{figure}

\subsection*{Dejavniki od tedna 45 do tedna 47}

\begin{table}[H]
\centering\small

\footnotesize\setlength{\tabcolsep}{4pt}
\begin{tabular}{lrrrl}
\toprule
dejavnik & teden 45 & teden 46 & teden 47 & vpliv na ceno t47 proti t45 \\
\midrule
veter v severni Nemčiji (km/h) & 17,4 & 23,2 & \textbf{29,5} & \textbf{močno navzdol} (poceni regija) \\
temperatura HU (°C) & 5,3 & 6,5 & 7,2 & navzdol: manj ogrevanja \\
cena plina TTF (€/MWh) & 90,2 & 87,4 & 91,0 & nevtralno \\
Paks (MW) & 1422 & 1458 & 1460 & nevtralno \\
sonce HU (MW) & 653 & 288 & 426 & navzgor: manj sonca kot v t45 \\
preostala poraba HU (MW) & 3349 & 3943 & 3581 & rahlo navzgor \\
cena HU (€/MWh) & 291,0 & 285,9 & ? & t46 je že pod t45 \\
cene AT / SK / DE (€/MWh) & 287 / 293 / 257 & 280 / 286 / 245 & ? & v t46 povsod nižje \\
\bottomrule
\end{tabular}
\caption{Dejavniki cene (tedenska povprečja; teden 47 = napoved z dne 14.\,11. ali zadnja znana vrednost).}
\end{table}

Domači dejavniki so mešani: toplejše vreme ceno znižuje, manj sonca kot v tednu 45 in nekoliko več preostale
porabe pa jo dvigujeta. Cena plina je praktično enaka, zato plin smeri ne določa. Odloči \textbf{veter v
Nemčiji}, ki bo po napovedi skoraj dvakrat močnejši kot v tednu 45. Ker je Madžarska tesno povezana s sosedi,
poceni elektrika v Nemčiji in Avstriji zniža ceno tudi na Madžarskem – pri sosedih so se cene v tednu 46 že
znižale.

\subsection*{Preverjanje s cenovnimi modeli}

Štiri cenovne modele smo preverili na 150 tednih (od 2024, ker so podatki o vetru v Nemčiji na voljo od 2023):
kolikokrat pravilno ugotovijo, ali bo teden dražji ali cenejši od tedna dva tedna prej (tako kot teden 47
proti tednu 45).

\begin{table}[H]
\centering\small

\footnotesize\setlength{\tabcolsep}{4pt}
\begin{tabular}{lrrr}
\toprule
model & napaka tedna (€/MWh) & pravilna smer & teden 47 (€/MWh) \\
\midrule
linearna regresija (TTF, preostala poraba, veter DE, dan) & \textbf{15,6} & 74\,\% & \textbf{263} \\
nevronska mreža (isti vhodi) & 16,8 & 75\,\% & 263 \\
naivni (kot prejšnji teden) & 17,3 & 74\,\% & 286 \\
merit order, le podatki HU (TTF × faktor) & 20,8 & 68\,\% & 293 \\
\bottomrule
\end{tabular}
\caption{Cenovni modeli: test 150 tednov in napoved za teden 47.}
\end{table}

Trije od štirih modelov in v testu najboljši model napovedujejo nižjo ceno. Merit order samo z madžarskimi
podatki se v testu odreže najslabše, ker madžarsko ceno močno določajo sosedje.

\textbf{Odločitev: povprečna cena v tednu 47 bo nižja od 290,98\,€/MWh}, približno 260–275\,€/MWh (base).
Povezava z b): isti vzrok – poceni elektrika v regiji ob vetru – pomeni tudi, da bodo madžarske plinske
elektrarne delale manj in bo uvoz večji.

\textbf{Tveganja:} močan mraz ali brezvetrje (napoved vetra se lahko spremeni), nenaden skok cene plina,
nenačrtovan izpad v Paksu ali pri sosedih. Odločitev se obrne, če bi bil veter v Nemčiji šibkejši od tedna 45.

\end{document}
